\documentclass[../../main.tex]{subfiles}
\begin{document}

\paragraph{【新課程】}
平地に3本のテレビ塔がある．
ひとりの男がこの平地の異なる3地点$A$，$B$，$C$に立って，その先端を眺めたところ，
どの地点でもそのうちの二つの先端が重なって見えた．
このとき，$A$，$B$，$C$は一直線上になければならない．
この理由を述べよ．

\paragraph{【旧課程】}
平面上に，どの3点も正三角形をつくらない相異なる4点がある．
これらを結ぶ6本の線分のうち，3本の長さは1であり，他の3本の長さは$a$であるという．
$a$の値を求めよ．

\end{document}
